\documentclass[5p,times]{elsarticle}

\usepackage{amsmath,amssymb,amsfonts}
\usepackage{algorithmic}
\usepackage{graphicx}
\usepackage{textcomp}
\usepackage{xcolor}
\usepackage{algorithm}
\usepackage{booktabs}
\usepackage[colorlinks=true, urlcolor=black, linkcolor=black, citecolor=black]{hyperref}

\begin{document}

\begin{frontmatter}

\title{Mimicking Biological Memory Consolidation in Spiking Neural Networks}

%% Elsevier Author Format:
\author[inst1]{Vishesh Kumar Singh}
\ead{visheshk24@iitk.ac.in}

\affiliation[inst1]{organization={Indian Institute of Technology Kanpur},
            % addressline={Kanpur}, 
            state={Uttar Pradesh},
            country={India}}

\begin{abstract}
Spiking Neural Networks (SNNs) offer a promising paradigm for energy-efficient neuromorphic computing, yet they pose significant challenges in maintaining knowledge across sequential tasks. In this work, we explore a biologically grounded framework for maintaining memory via synaptic consolidation. We investigate how specific neural patterns responsible for prediction can be identified in an SNN by introducing a \textbf{Potentiation-Factor (P-Factor)}, and propose a method to preserve the memory of past training while learning a new task. By tracking the potentiation of individual neurons based on local firing activity and reward signals, we implement a dual-phase learning method: (1) \textbf{Synaptic Locking} of identified expert neurons and (2) \textbf{Homeostatic Reset} of redundant novice neurons. Evaluated on the \textbf{Split-MNIST benchmark}, this approach demonstrates that the P-factor successfully identifies key neurons, achieving a peak \textbf{83.46\% retention} on the previous task while maintaining 90.63\% accuracy on the new task. These results suggest that mimicking the structural dynamics of the brain offers a viable pathway to lifelong learning in SNNs without the need for computationally expensive global curvature information (e.g., Fisher Information) or buffers of past data.
To facilitate reproducibility, our source code is publicly available at \url{https://github.com/vishesh-kumar-singh/Mimicking-Biological-Memory-Consolidation-in-SNNs}.
\end{abstract}

\begin{keyword}
%% Keywords are separated by \sep in Elsevier
Spiking Neural Networks \sep Continual Learning \sep Neuromorphic Computing \sep Synaptic Consolidation \sep Catastrophic Forgetting
\end{keyword}

\end{frontmatter}

\section{Introduction}
Human brain possesses remarkable ability to continuously learn new skills without erasing previously established memories, leading to lifetime learning. This capacity relies on the \textit{stability-plasticity dilemma} \cite{Mermillod2013}: neural circuits must remain plastic enough to encode novel information (learning) yet stable enough to preserve the traces of past experiences (memory). In biological systems, this balance is not achieved through global optimization or external data buffers, but through \textbf{Synaptic Consolidation}.

Neuroscience suggests that repeated activations of a neural pathway strengthens the connections between those neurons via triggering a cascade of chemical "tagging" processes (e.g., via CaMKII or PKM$\zeta$ proteins) that structurally reinforce specific synapses \cite{Clopath2010, Frankland2005}, leading to a phenomenon known as Long-Term Potentiation. This is the primary phenomenon is responsible for memory retention in biological brain. These reinforced neural pathways known engrams, locks and prevent degradation of memory, and this leaves remaining novice synapses free to learn new tasks

\subsection{Spiking Neural Networks as a Biological Substrate}
Spiking Neural Networks (SNNs) represent the ``third generation'' of neural models \cite{Maass1997}, designed to emulate the discrete, event-driven processing of the brain. Unlike traditional Artificial Neural Networks (ANNs) that rely on continuous activations, SNNs communicate via binary spikes, offering extreme energy efficiency suitable for neuromorphic hardware \cite{Roy2019, Davies2018}. 

However, standard SNN training algorithms (such as surrogate gradient descent \cite{Neftci2019}) treat all synaptic weights as equally plastic throughout the network's lifecycle. Due to this while adapting to a new task, the important weights for previous tasks are altered leading to complete overwriting of prior knowledge, a phenomenon known as Catastrophic Forgetting \cite{McCloskey1989}. In terms of Split-MNIST benchmark \cite{Zenke2017}, this leads to almost 0.00\% of retention of prior knowledge.

\subsection{The Gap: Biological Plausibility vs. Engineering fixes}
To mitigate forgetting, the field of Continual Learning (CL) has largely relied on engineering solutions that diverge from biological reality:
\begin{enumerate}
    \item \textbf{Regularization (e.g., EWC \cite{Kirkpatrick2017}):} These methods constrain weight updates based on global loss curvature (Fisher Information). While mathematically rigorous, calculating second-order derivatives (Hessian) is computationally expensive particularly over edge hardware like neuromorphic chips and biologically implausible because synapses cannot access global network curvature.
    \item \textbf{Replay Buffers \cite{Rebuffi2017}:} This includes storing subsets of raw past data for rehearsal during new training. It works, but human memory is more robust and doesn't need rehearsal again and again of prior knowledge while learning new one. We aimed to simulate this memory retention.
    \item \textbf{Parameter Isolation:} This includes the selection of a few parameters and freezing that part of the network.
\end{enumerate}

\subsection{Our Hypothesis: P-Factor Consolidation}
In this work, we investigate a fundamental question: \textit{Can simple, local Hebbian traces be sufficient to organize stable memory engrams in an SNN without global supervision?}

We propose a biologically-grounded framework termed as P-Factor Consolidation. Rather than seeing memory consolidation as a global optimization problem, we tried to model it as a local maturation process. We introduce a synaptic state variable, the Potentiation Factor(P-factor), which acts as a digital analog to biological synaptic tags. This factor evolves completely via the local activity of neurons during training via the following method:
\begin{itemize}
    \item \textbf{Potentiation:} Synapses that repeatedly participate in successful firing events accumulate $P$, signaling their reliability. Scaling these weights based on P-factor results in more independence of the correct neurons for predictions.
    \item \textbf{Engram Locking:} At task boundaries, neurons with high $P$-values are identified as the "Engram" and structurally frozen. The memory of the training is stored in these engrams.
    \item \textbf{Homeostatic Reset:} Low-confidence "Novice" neurons are re-initialized, mimicking biological turnover or neurogenesis to provide fresh capacity for new tasks.
\end{itemize}

By simulating this dual-phase process of \textbf{Locking (Stability)} and \textbf{Resetting (Plasticity)}, we aim to demonstrate that bio-inspired, local plasticity rules can spontaneously segregate the network into "Expert" and "Novice" partitions. Our empirical results on Split-MNIST, achieving up to \textbf{83.46\% retention}, serves as validation that mimicking the structural dynamics of the brain is a viable, energy-efficient pathway for lifelong learning.

% --- Related Work ---
\section{Related Work}
The challenge of Lifelong Learning (LL) has been extensively studied in Artificial Neural Networks (ANNs), but its translation to the energy-constrained, discrete domain of Spiking Neural Networks (SNNs) remains an active frontier. Existing approaches generally fall into three categories: Regularization, Replay, and Parameter Isolation.

\subsection{Regularization-Based Methods} Regularization approaches aim to penalize changes to important weights during subsequent tasks. The most prominent example, Elastic Weight Consolidation (EWC) \cite{Kirkpatrick2017}, computes the Fisher Information Matrix to estimate the "importance" of each parameter to the previous task's loss. While mathematically rigorous, EWC requires computing the second-derivative (Hessian) of the global loss, a computationally expensive operation that violates the locality constraints of neuromorphic hardware. Similarly, Synaptic Intelligence (SI) \cite{Zenke2017} accumulates a local path integral of gradient contributions to estimate synaptic importance online. However, both EWC and SI rely on high-precision gradient calculations during backpropagation, which are difficult to implement in SNNs where gradients are surrogate approximations \cite{Neftci2019}. In contrast, our P-Factor method avoids gradient-based importance estimation entirely, relying instead on forward-pass firing statistics (LTP/LTD) to determine synaptic utility.

\subsection{Replay and Generative Models} Replay-based methods mitigate forgetting by interleaving old data with new training samples \cite{Rebuffi2017}. In the absence of stored data, Generative Replay [Shin et al., 2017] trains a separate Generative Adversarial Network (GAN) or Autoencoder to "dream" pseudo-samples of past tasks. While effective, this approach doubles the computational overhead (training two networks) and requires complex reconstruction losses that are non-trivial to implement in spike-based systems. Our approach achieves retention without the need for external buffers or auxiliary generative networks, relying solely on the structural stability of the primary network.

\subsection{Parameter Isolation and Pruning} Our work is structurally most similar to parameter isolation methods like PackNet \cite{Mallya2018}, which iteratively prune and freeze network weights to create task-specific sub-networks. However, a critical distinction lies in the \textit{selection metric}. PackNet and similar magnitude-based pruning methods assume that \textit{"large weights are important."} In SNNs, this assumption is often flawed; a synapse with a large weight may connect to a neuron that rarely spikes due to high firing thresholds, rendering the weight functionally irrelevant. Preserving such "silent but large" weights wastes valuable capacity. Our P-Factor Consolidation diverges from this engineering heuristic by adopting a functional utility metric. By tracking the \textit{Potentiation Factor}, a history of successful firing events, we identify and freeze neurons that are actively participating in correct predictions, regardless of their raw weight magnitude. This ensures that the protected "engram" is composed of functionally active circuits rather than just statistically large parameters.

\subsection{Relation to Three-Factor Rules and E-prop}
Our P-Factor mechanism aligns with the \textbf{Three-Factor Learning Rule} framework \cite{Fremaux2016}, where plasticity is gated by a global neuromodulatory signal (Reward/Error). In modern SNN training, algorithms like \textbf{e-prop} \cite{Bellec2020} utilize similar local eligibility traces to approximate Backpropagation Through Time (BPTT). However, while e-prop solves the \textit{temporal credit assignment} problem for a single task, it does not inherently prevent catastrophic forgetting across sequential tasks. Our approach is distinct but complementary: whereas e-prop calculates the \textit{direction} of weight updates, P-Factor consolidation acts as a \textit{meta-plasticity gate}, determining the \textit{structural stability} of a neuron based on historical utility.

\subsection{Biological Models of Consolidation} In the domain of computational neuroscience, models like the Cascade Model \cite{Fusi2005} and Synaptic Tagging and Capture (STC) \cite{Clopath2008} have explored how synapses transition between plastic and stable states. Recently, approaches like \cite{Luboeinski2021} have proposed two-stage plasticity rules for burst-dependent consolidation. Our work bridges the gap between these theoretical neuroscience models and practical machine learning tasks. While Cascade and STC focus on single-synapse dynamics, we scale this principle to network-level task retention (Split-MNIST), demonstrating that local consolidation rules can effectively support supervised continual learning in deep SNN architectures.

\section{Methodology}

The core of our approach is the \textbf{Potentiation Factor ($P$)}, a local memory trace representing the structural stability and functional importance of a neuron. Unlike global importance metrics (e.g., Fisher Information) that require computing the Hessian over the entire loss landscape, our method relies on activity-dependent traces governed by local plasticity rules, making it highly efficient for neuromorphic implementations.

\subsection{Neuron-Centric P-Factor Dynamics}
In biological systems, synaptic tagging occurs at the level of individual dendritic spines (synapse-specific). However, implementing synapse-specific variables in SNNs scales with $\mathcal{O}(N^2)$ (number of connections), which is prohibitively memory-intensive for large-scale chips.

To address this, we introduce a \textbf{Neuron-Centric} approximation. We track the reliability of the \textit{post-synaptic neuron} as a proxy for the stability of its incoming weights. This reduces the memory overhead to $\mathcal{O}(N)$. 
Let $\mathbf{W} \in \mathbb{R}^{N_{out} \times N_{in}}$ be the weight matrix of the hidden layer, where the row vector $\mathbf{W}_{i,:}$ represents the incoming weights for neuron $i$. We introduce a state vector $\mathbf{P} \in \mathbb{R}^{N_{out}}$, initialized to zeros. The effective weights $\mathbf{W}_{eff}$ used during the forward pass are modulated by this factor:

\begin{equation}
    \mathbf{W}_{eff}[i, :] = \mathbf{W}[i, :] \cdot (1 + P_i)
\end{equation}

Biologically, this mimics the modulation of intrinsic excitability. By amplifying the effective synaptic strength of high-$P$ neurons, we create a "winner-take-all" dynamic during training. This ensures that "Expert" neurons respond more robustly to familiar patterns, naturally segregating them from "Novice" neurons, which remain quieter and plastic.

\subsection{Three-Factor Plasticity Rule}
The evolution of $P_i$ follows a biologically plausible \textbf{Three-Factor Rule}: (1) Pre-synaptic activity (Input), (2) Post-synaptic firing rate ($R_i$), and (3) A global neuromodulatory signal (Reward or Error). This structure bridges the gap between local Hebbian traces and the supervised nature of the task.

\subsubsection{Long-Term Potentiation (Reward)}
When the network correctly classifies an input, a global "Reward" signal triggers Long-Term Potentiation (LTP). Active neurons increase their P-factor to consolidate the successful circuit:
\begin{equation}
    P_i \leftarrow P_i + \alpha_{LTP} \cdot R_i \cdot (P_{max} - P_i)
\end{equation}
Here, $R_i$ is the firing rate (spike count) of neuron $i$. The term $(P_{max} - P_i)$ acts as a soft saturation boundary (negative feedback), preventing unbounded growth. This ensures that only neurons consistently involved in correct predictions reach the maximum expert status ($P \approx P_{max}$).


\subsubsection{Long-Term Depression (Punishment)}
Conversely, if the network makes an \textbf{incorrect prediction}, a global "Punishment" signal triggers Long-Term Depression (LTD). Active neurons (which contributed to the noise/error) are penalized:
\begin{equation}
    P_i \leftarrow \max(P_i - \alpha_{LTD} \cdot R_i, P_{min})
\end{equation}
We employ a linear decay for LTD to swiftly suppress erroneous neurons to $P_{min}$, effectively pruning unstable representations from the candidate engram.

Using Long-Term Potentiation and Long-Term Depression, we encourage the activity of selected correct neurons and discourage the activity of other neurons. This helps our engrams act independently of other nonessential neurons.

\subsection{Consolidation: Simulation of "Sleep"}
At the boundary between tasks, we simulate a period of offline consolidation. The accumulated P-factors are used to structurally reorganize the network into stable and plastic partitions:

\begin{enumerate}
    \item \textbf{Engram Identification:} We analyze the distribution of $P$ values across the hidden layer. A percentile threshold $\tau$ (e.g., 80th percentile) serves as the boundary. Neurons with $P_i > \text{Percentile}(\mathbf{P}, \tau)$ are identified as the \textbf{Engram}. These engrams contain the maximum information about the network or store the memory for the task it was trained on.
    
    \item \textbf{Synaptic Locking (Stability):} For every identified Engram neuron $i$, we freeze its \textbf{incoming synaptic weights} $\mathbf{W}_{i,:}$ (setting $\nabla \mathbf{W}_{i,:} = 0$). This ensures the neuron's feature selectivity remains invariant for future tasks and the memory is permanently locked in the neurons. Additionally, the output layer weights corresponding to Task A classes ($0-4$) are frozen to preserve the classifier mapping.
    
    \item \textbf{Homeostatic Reset (Plasticity):} Hidden neurons with low $P$ values are deemed redundant. Their incoming weights are re-initialized from a uniform distribution:
    % \begin{equation}
    %     \mathbf{W}_{novice} \sim \mathcal{U}(-\sqrt{k}, \sqrt{k})
    % \end{equation}
    \textit{Justification:} During Task A training, these "Novice" neurons were likely suppressed via LTD (treated as noise). Resetting them removes this "negative bias," providing a fresh, unbiased capacity for Task B, analogous to the biological phenomenon of neurogenesis. Similarly, the output weights for Task B classes ($5-9$) are reset to facilitate new learning.
\end{enumerate}

\section{Algorithm}

\begin{algorithm}[H] % IEEE prefers H over H
\caption{P-Factor Consolidation for Continual Learning}
\begin{algorithmic}[1]
\STATE \textbf{Input:} Model $\mathcal{M}$, Task A Data $\mathcal{D}_A$, Task B Data $\mathcal{D}_B$, Threshold $\tau$
\STATE \textbf{Initialize:} $P_i \leftarrow 0$ for all neurons

\STATE \textit{// Phase 1: Train Task A (Memory Formation)}
\FOR{epoch $e \in 1 \dots E_A$}
    \FOR{batch $(x, y) \in \mathcal{D}_A$}
        \STATE Forward pass: $\hat{y} = \mathcal{M}(x)$
        \STATE Update Weights $W$ via Backprop
        \IF{$\hat{y} == y$}
            \STATE $P_i \leftarrow P_i + \alpha_{LTP}(P_{max} - P_i)$ \textbf{for firing neurons}
        \ELSE
            \STATE $P_i \leftarrow P_i - \alpha_{LTD}$ \textbf{for firing neurons}
        \ENDIF
    \ENDFOR
\ENDFOR

\STATE \textit{// Phase 2: Consolidation (Sleep)}
\STATE Calculate Threshold $\theta = \text{Percentile}(P, \tau)$
\STATE $\text{Mask}_{expert} \leftarrow (P \geq \theta)$
\STATE $\text{Mask}_{novice} \leftarrow (P < \theta)$
\STATE \textbf{Freeze:} Lock weights where $\text{Mask}_{expert}$ is True
\STATE \textbf{Reset:} Re-init weights where $\text{Mask}_{novice}$ is True
\STATE \textbf{Liberate:} Reset $P$ and $W$ for Task B Output Heads

\STATE \textit{// Phase 3: Train Task B (New Learning)}
\FOR{epoch $e \in 1 \dots E_B$}
    \FOR{batch $(x, y) \in \mathcal{D}_B$}
        \STATE Forward pass: $\hat{y} = \mathcal{M}(x)$
        \STATE Compute Gradients $\nabla W$
        \STATE $\nabla W[\text{Mask}_{expert}] \leftarrow 0$ \COMMENT{Apply Mask}
        \STATE Update $W$
    \ENDFOR
\ENDFOR
\end{algorithmic}
\end{algorithm}

\section{Experimental Setup}

\subsection{Datasets and Task Formulation}
% [Placeholder: Define Split-MNIST and N-MNIST (Classes 0-4 vs 5-9)]
% [Placeholder: Define single-head, Class-Incremental Learning (CIL) vs Task-Incremental Learning (TIL) setup]

\subsection{Network Architecture}
% [Placeholder: Define 784/N-MNIST-input -> 1024 -> 10 topology]
% [Placeholder: Define LIF neuron dynamics and surrogate gradients]

\subsection{Implementation Details}
% [Placeholder: Define Optimizer (Adam), Learning Rate, Batch Size]
% [Placeholder: Define P-Factor hyperparameters (alpha_LTP, alpha_LTD, thresholds)]

\subsection{Evaluation Metrics}
% [Placeholder: Define Task A Retention (Stability) and Task B Accuracy (Plasticity)]
% [Placeholder: Define Combined Score]

\subsection{Training Protocol and Baselines}
% [Placeholder: Describe Phase 1 (Maturation) and Phase 2 (Adaptation) across Epochs \in {1, 3, 5}]
% [Placeholder: List Baselines: Fine-tuning, ER, EWC, SI, PackNet, Random, Index]


\section{Results and Discussion}
% [Placeholder: Brief statement that all results are averaged over 5 seeds (42-46)]

\subsection{The Failure of Gradient-Based Regularization in SNNs}
% [Placeholder: Text introducing the failure of EWC (FIM sparsity) and SI (Surrogate gradient noise)]

\begin{table}[H]
    \caption{Gradient-Based Baselines Performance Summary (Split-MNIST, Epoch 3)}
    \begin{center}
    \resizebox{\columnwidth}{!}{
    \begin{tabular}{lccc}
        \toprule
        \textbf{Method} & \textbf{Task A (Ret.)} & \textbf{Task B (Acc.)} & \textbf{Combined} \\
        \midrule
        Fine-Tuning (Baseline) & -- & -- & -- \\
        \midrule
        EWC ($\lambda=10^3$) & -- & -- & -- \\
        EWC ($\lambda=10^5$) & -- & -- & -- \\
        EWC ($\lambda=10^6$) & -- & -- & -- \\
        \midrule
        SI ($c=1$) & -- & -- & -- \\
        SI ($c=10^2$) & -- & -- & -- \\
        SI ($c=10^4$) & -- & -- & -- \\
        SI ($c=10^6$) & -- & -- & -- \\
        \bottomrule
    \end{tabular}}
    \label{tab:gradient_baselines}
    \end{center}
\end{table}

\subsection{Structural Consolidation: Deterministic vs. Stochastic}
% [Placeholder: Text comparing P-Factor to PackNet, Random, and Fixed Index]

\begin{table}[H]
    \caption{Comparison of Structural Engram Allocation Strategies (Split-MNIST, Epoch 3)}
    \begin{center}
    \resizebox{\columnwidth}{!}{
    \begin{tabular}{lccc}
        \toprule
        \textbf{Configuration} & \textbf{Task A (Ret.)} & \textbf{Task B (Acc.)} & \textbf{Combined} \\
        \midrule
        \multicolumn{4}{l}{\textit{Threshold $\tau=0.2$ (20\% Frozen)}} \\
        PackNet (Magnitude) & -- & -- & -- \\
        Random Freezing & -- & -- & -- \\
        Fixed Index & -- & -- & -- \\
        \textbf{P-Factor (Ours)} & -- & -- & -- \\
        \midrule
        \multicolumn{4}{l}{\textit{Threshold $\tau=0.4$ (40\% Frozen)}} \\
        PackNet (Magnitude) & -- & -- & -- \\
        Random Freezing & -- & -- & -- \\
        Fixed Index & -- & -- & -- \\
        \textbf{P-Factor (Ours)} & -- & -- & -- \\
        \midrule
        \multicolumn{4}{l}{\textit{Threshold $\tau=0.6$ (60\% Frozen)}} \\
        PackNet (Magnitude) & -- & -- & -- \\
        Random Freezing & -- & -- & -- \\
        Fixed Index & -- & -- & -- \\
        \textbf{P-Factor (Ours)} & -- & -- & -- \\
        \midrule
        \multicolumn{4}{l}{\textit{Threshold $\tau=0.8$ (80\% Frozen)}} \\
        PackNet (Magnitude) & -- & -- & -- \\
        Random Freezing & -- & -- & -- \\
        Fixed Index & -- & -- & -- \\
        \textbf{P-Factor (Ours)} & -- & -- & -- \\
        \bottomrule
    \end{tabular}}
    \label{tab:structural_comparison}
    \end{center}
\end{table}

\subsection{Longitudinal Dynamics: The Stability-Plasticity Trajectory}
% [Placeholder: Text discussing Phase 1 (Maturation), Phase 2 (Peak), Phase 3 (Plasticity Creep)]

\begin{figure}[H]
    % \centerline{\includegraphics[width=\columnwidth]{longitudinal_trajectory.png}}
    \caption{Average stability and plasticity trajectories across Epochs 1, 3, and 5 for $\tau=0.8$.}
    \label{fig:longitudinal_trajectory}
\end{figure}

\subsection{Ablation Studies: Mechanics of the Engram}
% [Placeholder: Text introducing ablations on scaling and reset strategies]

\begin{table}[H]
    \caption{Ablation Study: Forward Pass Scaling and Reset Mechanics (Split-MNIST, Epoch 3)}
    \begin{center}
    \resizebox{\columnwidth}{!}{
    \begin{tabular}{lccc}
        \toprule
        \textbf{Configuration} & \textbf{Task A (Ret.)} & \textbf{Task B (Acc.)} & \textbf{Combined} \\
        \midrule
        \multicolumn{4}{l}{\textit{Weight Scaling Ablation ($\tau=0.8$)}} \\
        No-Scale & -- & -- & -- \\
        \textbf{P-Factor (With Scale)} & -- & -- & -- \\
        \midrule
        \multicolumn{4}{l}{\textit{Novice Reset Strategy ($\tau=0.8$)}} \\
        No-Reset & -- & -- & -- \\
        Reset-Zero & -- & -- & -- \\
        Reset-Scale & -- & -- & -- \\
        \textbf{Kaiming Reset (Ours)} & -- & -- & -- \\
        \bottomrule
    \end{tabular}}
    \label{tab:ablations}
    \end{center}
\end{table}

\subsection{Scaling to High-Entropy Temporal Streams (N-MNIST)}
% [Placeholder: Text discussing Class-IL collapse and the necessity of manifold coverage via stochastic freezing in temporal domains]

\begin{table}[H]
    \caption{Performance on Spatio-Temporal Event Streams (N-MNIST, Epoch 3)}
    \begin{center}
    \resizebox{\columnwidth}{!}{
    \begin{tabular}{lccc}
        \toprule
        \textbf{Configuration} & \textbf{Task A (Ret.)} & \textbf{Task B (Acc.)} & \textbf{Combined} \\
        \midrule
        Fine-Tuning (Baseline) & -- & -- & -- \\
        ER (Buffer=200) [Oracle] & -- & -- & -- \\
        \midrule
        Random Freezing ($\tau=0.4$) & -- & -- & -- \\
        Random Freezing ($\tau=0.6$) & -- & -- & -- \\
        Random Freezing ($\tau=0.8$) & -- & -- & -- \\
        \midrule
        \textbf{P-Factor ($\tau=0.4$)} & -- & -- & -- \\
        \textbf{P-Factor ($\tau=0.6$)} & -- & -- & -- \\
        \textbf{P-Factor ($\tau=0.8$)} & -- & -- & -- \\
        \bottomrule
    \end{tabular}}
    \label{tab:nmnist_results}
    \end{center}
\end{table}

\subsection{Energy and Computational Cost Analysis}
% [Placeholder: Text discussing the O(N) vs O(W) memory overhead and energy efficiency extracted from analyze_energy.py]

\begin{table}[H]
    \caption{Computational Overhead and Energy Efficiency Comparison}
    \begin{center}
    \resizebox{\columnwidth}{!}{
    \begin{tabular}{lccc}
        \toprule
        \textbf{Method} & \textbf{Memory Overhead} & \textbf{Relative Energy Cost} & \textbf{Complexity} \\
        \midrule
        Fine-Tuning & None & 1.0x & $O(1)$ \\
        ER (Buffer=200) & High (Raw Data) & -- & $O(D_{mem})$ \\
        EWC & High (Float Matrix) & -- & $O(W)$ \\
        SI & High (Float Matrix) & -- & $O(W)$ \\
        \textbf{P-Factor} & \textbf{Low (Binary Mask)} & \textbf{--} & \textbf{$O(N)$} \\
        \bottomrule
    \end{tabular}}
    \label{tab:energy_analysis}
    \end{center}
\end{table}

\section{Conclusion}
In this work, we presented P-Factor Consolidation, a novel framework that bridges the gap between biological synaptic tagging and practical continual learning in Spiking Neural Networks. By introducing a local, activity-dependent state variable, we demonstrated that the complex phenomenon of memory consolidation can be effectively mimicked without the need for computationally expensive global curvature estimates (e.g., Hessian matrices) or external replay buffers.

Our empirical results on the Split-MNIST benchmark validate that the P-Factor mechanism successfully identifies and isolates a functional "engram", a sparse sub-network responsible for past tasks, achieving a peak retention of \textbf{83.46\%}. Crucially, the segregation of "Expert" and "Novice" neurons via the Three-Factor Rule proves that structural plasticity alone, governed by local firing history, is sufficient to mitigate catastrophic forgetting.

While the current implementation relies on a fixed consolidation threshold ($\tau=0.8$), which imposes a finite capacity limit for very long task sequences, it establishes a foundational proof-of-concept for autonomous neuromorphic agents. Future work will explore dynamic thresholding mechanisms and synaptic regrowth strategies to further extend the network's plasticity, moving closer to truly open-ended lifelong learning on resource-constrained edge hardware.

\section*{Declaration of generative AI and AI-assisted technologies in the manuscript preparation process.}
During the preparation of this work, the author used Gemini in order to refine the content of the manuscript. After using this, the author reviewed and edited the content as needed and takes full responsibility for the content of the published article.

% \clearpage
\begin{thebibliography}{00}

\bibitem{Mermillod2013} M. Mermillod, A. Bugaiska, and P. Bonin, ``The stability-plasticity dilemma: Investigating the continuum from catastrophic forgetting to age-limited learning effects,'' \textit{Frontiers in Psychology}, vol. 4, p. 504, 2013.

\bibitem{Clopath2010} C. Clopath, L. B\"using, E. Vasilaki, and W. Gerstner, ``Connectivity reflects coding: a model of voltage-based STDP with homeostasis,'' \textit{Nature Neuroscience}, vol. 13, no. 3, pp. 344--352, 2010.

\bibitem{Frankland2005} P. W. Frankland and B. Bontempi, ``The organization of recent and remote memories,'' \textit{Nature Reviews Neuroscience}, vol. 6, no. 2, pp. 119--130, 2005.

\bibitem{Maass1997} W. Maass, ``Networks of spiking neurons: the third generation of neural network models,'' \textit{Neural Networks}, vol. 10, no. 9, pp. 1659--1671, 1997.

\bibitem{Roy2019} K. Roy, A. Jaiswal, and P. Panda, ``Towards spike-based machine intelligence with neuromorphic computing,'' \textit{Nature}, vol. 575, no. 7784, pp. 607--617, 2019.

\bibitem{Davies2018} M. Davies et al., ``Loihi: A neuromorphic manycore processor with on-chip learning,'' \textit{IEEE Micro}, vol. 38, no. 1, pp. 82--99, 2018.

\bibitem{Neftci2019} E. O. Neftci, H. Mostafa, and F. Zenke, ``Surrogate gradient learning in spiking neural networks,'' \textit{IEEE Signal Processing Magazine}, vol. 36, no. 6, pp. 51--63, 2019.

\bibitem{McCloskey1989} M. McCloskey and N. J. Cohen, ``Catastrophic interference in connectionist networks: The sequential learning problem,'' \textit{Psychology of Learning and Motivation}, vol. 24, pp. 109--165, 1989.

\bibitem{Zenke2017} F. Zenke, B. Poole, and S. Ganguli, ``Continual learning through synaptic intelligence,'' in \textit{Proc. Int. Conf. Mach. Learn. (ICML)}, 2017, pp. 3987--3995.

\bibitem{Kirkpatrick2017} J. Kirkpatrick et al., ``Overcoming catastrophic forgetting in neural networks,'' \textit{Proc. Natl. Acad. Sci. (PNAS)}, vol. 114, no. 13, pp. 3521--3526, 2017.

\bibitem{Rebuffi2017} S.-A. Rebuffi, A. Kolesnikov, G. Sperl, and C. H. Lampert, ``iCaRL: Incremental classifier and representation learning,'' in \textit{Proc. IEEE Conf. Comput. Vis. Pattern Recognit. (CVPR)}, 2017, pp. 2001--2010.

\bibitem{Mallya2018} A. Mallya and S. Lazebnik, ``PackNet: Adding multiple tasks to a single network by iterative pruning,'' in \textit{Proc. IEEE Conf. Comput. Vis. Pattern Recognit. (CVPR)}, 2018, pp. 7765--7773.

\bibitem{Fusi2005} S. Fusi, P. J. Drew, and L. F. Abbott, ``Cascade models of synaptically stored memories,'' \textit{Neuron}, vol. 45, no. 4, pp. 599--611, 2005.

\bibitem{Clopath2008} C. Clopath, L. Ziegler, E. Vasilaki, L. B\"using, and W. Gerstner, ``Tag-trigger-consolidation: a model of early and late long-term-potentiation and depression,'' \textit{PLoS Comput. Biol.}, vol. 4, no. 1, p. e1000002, 2008.

\bibitem{Luboeinski2021} J. Luboeinski and C. Tetzlaff, ``Memory consolidation and improvement by synaptic tagging and capture in recurrent neural networks,'' \textit{Commun. Biol.}, vol. 4, no. 1, p. 275, 2021.

\bibitem{Fremaux2016} N. Frémaux and W. Gerstner, ``Neuromodulated spike-timing-dependent plasticity, and theory of three-factor learning rules,'' \textit{Frontiers in Neural Circuits}, vol. 9, p. 85, 2016.

\bibitem{Bellec2020} G. Bellec et al., ``A solution to the learning dilemma for recurrent networks of spiking neurons,'' \textit{Nature Communications}, vol. 11, no. 1, p. 3625, 2020.

\end{thebibliography}


\appendix

\section{Detailed Experimental Results}
For completeness, we provide the full breakdown of experimental results across all training durations ($E \in \{1, 2, 3, 4, 5\}$) and consolidation thresholds ($\tau \in [0.4, 0.8]$).

% ================= EPOCH 1 =================
\subsection{Epoch 1 Results}

\begin{table}[H]
    \caption{\textbf{Epoch 1: P-Factor vs Baselines (Mean $\pm$ Std)}.}
    \begin{center}
    \resizebox{\columnwidth}{!}{
    \begin{tabular}{lccc}
        \toprule
        \textbf{Method} & \textbf{Task A (Ret.)} & \textbf{Task B (Acc.)} & \textbf{Combined} \\
        \midrule
        Fine-Tuning (Baseline) & 0.04\% $\pm$ 0.07 & 95.59\% $\pm$ 1.51 & 46.85\% \\
        ER (Buffer=200) & 89.46\% $\pm$ 2.58 & 95.29\% $\pm$ 0.79 & 92.65\% \\
        \midrule
        EWC ($\lambda=10^3$) & 2.13\% $\pm$ 3.57 & 96.15\% $\pm$ 0.64 & 48.40\% \\
        EWC ($\lambda=10^5$) & 9.68\% $\pm$ 11.85 & 79.75\% $\pm$ 19.81 & 44.21\% \\
        EWC ($\lambda=10^6$) & 24.54\% $\pm$ 5.55 & 25.25\% $\pm$ 21.43 & 24.52\% \\
        \midrule
        SI ($c=1$) & 1.43\% $\pm$ 2.86 & 96.33\% $\pm$ 0.78 & 48.13\% \\
        SI ($c=10^2$) & 25.92\% $\pm$ 15.29 & 94.09\% $\pm$ 1.55 & 59.38\% \\
        SI ($c=10^4$) & 28.42\% $\pm$ 11.85 & 93.10\% $\pm$ 1.15 & 60.45\% \\
        SI ($c=10^6$) & 56.67\% $\pm$ 11.27 & 87.66\% $\pm$ 6.75 & 72.18\% \\
        \midrule
        PackNet (20\%) & 0.22\% $\pm$ 0.32 & 95.23\% $\pm$ 1.51 & 46.99\% \\
        PackNet (40\%) & 2.12\% $\pm$ 1.47 & 95.83\% $\pm$ 0.66 & 48.27\% \\
        PackNet (60\%) & 9.69\% $\pm$ 1.65 & 96.39\% $\pm$ 0.78 & 52.36\% \\
        PackNet (80\%) & 31.39\% $\pm$ 8.71 & 96.33\% $\pm$ 0.34 & 63.57\% \\
        \midrule
        Random (20\%) & 1.35\% $\pm$ 1.54 & 95.72\% $\pm$ 1.02 & 47.67\% \\
        Random (40\%) & 9.99\% $\pm$ 3.33 & 96.05\% $\pm$ 0.80 & 52.33\% \\
        Random (60\%) & 20.27\% $\pm$ 8.27 & 95.29\% $\pm$ 1.62 & 57.22\% \\
        Random (80\%) & 39.25\% $\pm$ 4.84 & 94.54\% $\pm$ 1.31 & 66.62\% \\
        \midrule
        \textbf{P-Factor $\tau=0.2$ (Ours)} & -- & -- & -- \\
        \textbf{P-Factor $\tau=0.4$ (Ours)} & -- & -- & -- \\
        \textbf{P-Factor $\tau=0.6$ (Ours)} & -- & -- & -- \\
        \textbf{P-Factor $\tau=0.8$ (Ours)} & -- & -- & -- \\
        \bottomrule
    \end{tabular}}
    \end{center}
\end{table}

\begin{table}[H]
    \caption{\textbf{Epoch 1: Ablation Studies (Mean $\pm$ Std)}.}
    \begin{center}
    \resizebox{\columnwidth}{!}{
    \begin{tabular}{lccc}
        \toprule
        \textbf{Method} & \textbf{Task A (Ret.)} & \textbf{Task B (Acc.)} & \textbf{Combined} \\
        \midrule
        \multicolumn{4}{l}{\textit{Weight Scaling Ablation}} \\
        No-Scale $\tau=0.2$ & -- & -- & -- \\
        No-Scale $\tau=0.4$ & -- & -- & -- \\
        No-Scale $\tau=0.6$ & -- & -- & -- \\
        No-Scale $\tau=0.8$ & -- & -- & -- \\
        \midrule
        \multicolumn{4}{l}{\textit{Novice Reset Strategy}} \\
        No-Reset $\tau=0.2$ & -- & -- & -- \\
        No-Reset $\tau=0.4$ & -- & -- & -- \\
        No-Reset $\tau=0.6$ & -- & -- & -- \\
        No-Reset $\tau=0.8$ & -- & -- & -- \\
        \midrule
        Reset-Zero $\tau=0.2$ & -- & -- & -- \\
        Reset-Zero $\tau=0.4$ & -- & -- & -- \\
        Reset-Zero $\tau=0.6$ & -- & -- & -- \\
        Reset-Zero $\tau=0.8$ & -- & -- & -- \\
        \midrule
        Reset-Scale $\tau=0.2$ & -- & -- & -- \\
        Reset-Scale $\tau=0.4$ & -- & -- & -- \\
        Reset-Scale $\tau=0.6$ & -- & -- & -- \\
        Reset-Scale $\tau=0.8$ & -- & -- & -- \\
        \bottomrule
    \end{tabular}}
    \end{center}
\end{table}

% ================= EPOCH 3 =================
\subsection{Epoch 3 Results}

\begin{table}[H]
    \caption{\textbf{Epoch 3: P-Factor vs Baselines (Mean $\pm$ Std)}.}
    \begin{center}
    \resizebox{\columnwidth}{!}{
    \begin{tabular}{lccc}
        \toprule
        \textbf{Method} & \textbf{Task A (Ret.)} & \textbf{Task B (Acc.)} & \textbf{Combined} \\
        \midrule
        Fine-Tuning (Baseline) & 0.00\% $\pm$ 0.00 & 95.60\% $\pm$ 1.02 & 48.09\% \\
        ER (Buffer=200) & 90.70\% $\pm$ 1.23 & 95.38\% $\pm$ 0.70 & 94.24\% \\
        \midrule
        EWC ($\lambda=10^3$) & 0.59\% $\pm$ 0.93 & 95.51\% $\pm$ 0.60 & 48.25\% \\
        EWC ($\lambda=10^5$) & 15.89\% $\pm$ 7.48 & 34.10\% $\pm$ 36.46 & 32.51\% \\
        EWC ($\lambda=10^6$) & 23.62\% $\pm$ 3.63 & 4.11\% $\pm$ 8.21 & 15.54\% \\
        \midrule
        SI ($c=1$) & 7.60\% $\pm$ 4.22 & 96.11\% $\pm$ 0.99 & 51.85\% \\
        SI ($c=10^2$) & 18.55\% $\pm$ 8.29 & 94.53\% $\pm$ 0.91 & 56.49\% \\
        SI ($c=10^4$) & 27.30\% $\pm$ 6.73 & 81.96\% $\pm$ 14.39 & 60.92\% \\
        SI ($c=10^6$) & 54.97\% $\pm$ 11.02 & 66.47\% $\pm$ 17.01 & 68.18\% \\
        \midrule
        PackNet (20\%) & 1.02\% $\pm$ 0.87 & 95.85\% $\pm$ 1.04 & 48.57\% \\
        PackNet (40\%) & 8.99\% $\pm$ 1.15 & 95.51\% $\pm$ 1.39 & 52.79\% \\
        PackNet (60\%) & 16.22\% $\pm$ 2.37 & 96.00\% $\pm$ 0.75 & 56.52\% \\
        PackNet (80\%) & 34.27\% $\pm$ 5.06 & 95.63\% $\pm$ 0.93 & 65.86\% \\
        \midrule
        Random (20\%) & -- & -- & -- \\
        Random (40\%) & -- & -- & -- \\
        Random (60\%) & -- & -- & -- \\
        Random (80\%) & -- & -- & -- \\
        \midrule
        \textbf{P-Factor $\tau=0.2$ (Ours)} & -- & -- & -- \\
        \textbf{P-Factor $\tau=0.4$ (Ours)} & -- & -- & -- \\
        \textbf{P-Factor $\tau=0.6$ (Ours)} & -- & -- & -- \\
        \textbf{P-Factor $\tau=0.8$ (Ours)} & -- & -- & -- \\
        \bottomrule
    \end{tabular}}
    \end{center}
\end{table}

\begin{table}[H]
    \caption{\textbf{Epoch 3: Ablation Studies (Mean $\pm$ Std)}.}
    \begin{center}
    \resizebox{\columnwidth}{!}{
    \begin{tabular}{lccc}
        \toprule
        \textbf{Method} & \textbf{Task A (Ret.)} & \textbf{Task B (Acc.)} & \textbf{Combined} \\
        \midrule
        \multicolumn{4}{l}{\textit{Weight Scaling Ablation}} \\
        No-Scale $\tau=0.2$ & -- & -- & -- \\
        No-Scale $\tau=0.4$ & -- & -- & -- \\
        No-Scale $\tau=0.6$ & -- & -- & -- \\
        No-Scale $\tau=0.8$ & -- & -- & -- \\
        \midrule
        \multicolumn{4}{l}{\textit{Novice Reset Strategy}} \\
        No-Reset $\tau=0.2$ & -- & -- & -- \\
        No-Reset $\tau=0.4$ & -- & -- & -- \\
        No-Reset $\tau=0.6$ & -- & -- & -- \\
        No-Reset $\tau=0.8$ & -- & -- & -- \\
        \midrule
        Reset-Zero $\tau=0.2$ & -- & -- & -- \\
        Reset-Zero $\tau=0.4$ & -- & -- & -- \\
        Reset-Zero $\tau=0.6$ & -- & -- & -- \\
        Reset-Zero $\tau=0.8$ & -- & -- & -- \\
        \midrule
        Reset-Scale $\tau=0.2$ & -- & -- & -- \\
        Reset-Scale $\tau=0.4$ & -- & -- & -- \\
        Reset-Scale $\tau=0.6$ & -- & -- & -- \\
        Reset-Scale $\tau=0.8$ & -- & -- & -- \\
        \bottomrule
    \end{tabular}}
    \end{center}
\end{table}

% ================= EPOCH 5 =================
\subsection{Epoch 5 Results}

\begin{table}[H]
    \caption{\textbf{Epoch 5: P-Factor vs Baselines (Mean $\pm$ Std)}.}
    \begin{center}
    \resizebox{\columnwidth}{!}{
    \begin{tabular}{lccc}
        \toprule
        \textbf{Method} & \textbf{Task A (Ret.)} & \textbf{Task B (Acc.)} & \textbf{Combined} \\
        \midrule
        Fine-Tuning (Baseline) & 0.00\% $\pm$ 0.00 & 95.89\% $\pm$ 0.98 & 48.09\% \\
        ER (Buffer=200) & 90.93\% $\pm$ 0.93 & 93.93\% $\pm$ 1.31 & 94.42\% \\
        \midrule
        EWC ($\lambda=10^3$) & 3.71\% $\pm$ 4.59 & 95.75\% $\pm$ 1.21 & 50.14\% \\
        EWC ($\lambda=10^5$) & 14.79\% $\pm$ 0.00 & 71.91\% $\pm$ 0.00 & 56.05\% \\
        EWC ($\lambda=10^6$) & -- & -- & -- \\
        \midrule
        SI ($c=1$) & -- & -- & -- \\
        SI ($c=10^2$) & -- & -- & -- \\
        SI ($c=10^4$) & -- & -- & -- \\
        SI ($c=10^6$) & -- & -- & -- \\
        \midrule
        PackNet (20\%) & -- & -- & -- \\
        PackNet (40\%) & -- & -- & -- \\
        PackNet (60\%) & -- & -- & -- \\
        PackNet (80\%) & -- & -- & -- \\
        \midrule
        Random (20\%) & -- & -- & -- \\
        Random (40\%) & -- & -- & -- \\
        Random (60\%) & -- & -- & -- \\
        Random (80\%) & -- & -- & -- \\
        \midrule
        \textbf{P-Factor $\tau=0.2$ (Ours)} & -- & -- & -- \\
        \textbf{P-Factor $\tau=0.4$ (Ours)} & -- & -- & -- \\
        \textbf{P-Factor $\tau=0.6$ (Ours)} & -- & -- & -- \\
        \textbf{P-Factor $\tau=0.8$ (Ours)} & -- & -- & -- \\
        \bottomrule
    \end{tabular}}
    \end{center}
\end{table}

\begin{table}[H]
    \caption{\textbf{Epoch 5: Ablation Studies (Mean $\pm$ Std)}.}
    \begin{center}
    \resizebox{\columnwidth}{!}{
    \begin{tabular}{lccc}
        \toprule
        \textbf{Method} & \textbf{Task A (Ret.)} & \textbf{Task B (Acc.)} & \textbf{Combined} \\
        \midrule
        \multicolumn{4}{l}{\textit{Weight Scaling Ablation}} \\
        No-Scale $\tau=0.2$ & -- & -- & -- \\
        No-Scale $\tau=0.4$ & -- & -- & -- \\
        No-Scale $\tau=0.6$ & -- & -- & -- \\
        No-Scale $\tau=0.8$ & -- & -- & -- \\
        \midrule
        \multicolumn{4}{l}{\textit{Novice Reset Strategy}} \\
        No-Reset $\tau=0.2$ & -- & -- & -- \\
        No-Reset $\tau=0.4$ & -- & -- & -- \\
        No-Reset $\tau=0.6$ & -- & -- & -- \\
        No-Reset $\tau=0.8$ & -- & -- & -- \\
        \midrule
        Reset-Zero $\tau=0.2$ & -- & -- & -- \\
        Reset-Zero $\tau=0.4$ & -- & -- & -- \\
        Reset-Zero $\tau=0.6$ & -- & -- & -- \\
        Reset-Zero $\tau=0.8$ & -- & -- & -- \\
        \midrule
        Reset-Scale $\tau=0.2$ & -- & -- & -- \\
        Reset-Scale $\tau=0.4$ & -- & -- & -- \\
        Reset-Scale $\tau=0.6$ & -- & -- & -- \\
        Reset-Scale $\tau=0.8$ & -- & -- & -- \\
        \bottomrule
    \end{tabular}}
    \end{center}
\end{table}

\end{document}